% ============================================
% 第3章："这个学了有什么用？"——效率如何杀死了意义
% 梁桥 · 效率陷阱
% ============================================

\chapter{``这个学了有什么用？''——效率如何杀死了意义}

\vspace{0.5cm}

\begin{center}
{\small\color{muted} 梁桥是最简单的。但梁桥有一个致命的缺陷：}
{\small\color{muted} 跨度越大，自重越大。80\%的自重，20\%的荷载。}
{\small\color{muted} 效率极高。意义极低。}
\end{center}

\vspace{1cm}

\section{那个问题}

每一个教土木的老师，都被问过这句话。

可能在课堂上，一个坐在后排的学生，没有举手，直接问了出来。可能在办公室答疑，学生犹豫了很久，最后小声说了一句。可能在毕业聚餐上，几杯酒之后，终于说了实话。

\begin{center}
{\large\itshape ``老师，学这个，到底有什么用？''}
\end{center}

这个问题，我听过很多次。每一次，我都回答。我说，桥梁工程是基础设施的基石。我说，中国有超过一百万座公路桥，每年还需要新建几万座。我说，毕业生的就业率不低，薪资也不差。

但我知道，这些回答，没有一个真正回答了那个问题。

因为学生问的不是``学这个能找到什么工作''。学生问的是：\textbf{在AI能做一切的时代，我花四年时间学一门手艺，这件事本身，还值得吗？}

\section{``为了这碟醋包了顿饺子''}

2025年10月21日，我在课堂上讲桥梁的载重比。我讲了一段话，后来被转写成了文字：

\begin{shiyu-inline}
``我们的桥梁的载重比其实是非常小的。一般来说，一个大跨度桥梁，大概百分之八十是它的自重，百分之二十是车辆。所以说，我们实际上只有百分之二十的承载能力是用在了承载荷载上。所以这个就是很大的浪费嘛。就是姜文的电影是吧——\textbf{为了这碟醋包了顿饺子}。''

\hfill ——2025年10月21日，桥梁课程，录音转写
\end{shiyu-inline}

我当时说这句话的时候，教室里笑了一下。姜文的电影梗，大家都懂。但后来我重新读到这段转写，发现这个比喻比我想象的更深。

\textbf{梁桥80\%的自重，就是当代教育80\%的``无用功''。}

我们常常把大量时间用于公式复述、规范检索、软件操作与考试组织。AI正在加速其中一部分工作，于是那个旧问题变得更响：省下来的时间，能否真正还给尝试、解释、核验和创造？这里不再套用一个没有调查依据的百分比。

我们只有20\%的精力用在真正重要的东西上——\textbf{困惑、尝试、失败、骄傲、判断、审美、责任}。

就像梁桥，80\%的自重，20\%的荷载。\textbf{为了这碟醋，包了顿饺子。}

\section{载重比：一个工程概念如何变成教育隐喻}

在那堂课上，我还讲了另一个故事——关于结构设计大赛：

\begin{shiyu-inline}
``我们做了一个100克的船，它可以承载19公斤的重量。大概载重比是190倍。但是我们的桥梁，大跨度桥梁，载重比是非常小的。一百米的桥，大概百分之八十是它的自重。''

\hfill ——2025年10月21日，同一堂课
\end{shiyu-inline}

\textbf{载重比}，是一个工程概念。它衡量的是：一座桥的重量里，有多少是``必须的负担''（自重），有多少是``真正的价值''（承载交通）。

100克的船，承载19公斤——载重比190倍。100米的桥，80\%是自重——载重比1.25倍。

\textbf{载重比，也是一个教育概念。}

一门课的载重比是多少？学生花了80个小时，有多少小时是在``真正理解''，有多少小时是在``完成要求''？老师花了80个小时备课，有多少小时是在``设计学习体验''，有多少小时是在``做PPT排版''？

梁桥的困境，就是教育的困境：\textbf{我们花了太多精力做那些``必须做但没什么用''的事，以至于没时间做``真正重要''的事。}

\section{AI让这个问题变得更尖锐了}

2026年6月，我在知识星球上看到一段话，是师宇写的：

\begin{shiyu-inline}
``AI时代的教育，不该再假装学生生活在校园平台里。很多教育AI产品喜欢做封闭知识库、智能问答、自测系统。它们假设学生愿意在学校提供的环境里耐心学习。但真实世界不是这样运行的。学生未来面对的是开放工具、真实任务、复杂工作流和不断变化的平台。''

\hfill ——师宇博文，2026年6月8日
\end{shiyu-inline}

AI把这个问题逼到了极致。因为AI不只让``知识传递''变快了——它让``知识传递''\textbf{变得几乎免费}。

你不需要花四年时间学``桥梁工程基础知识''。你只需要打开ChatGPT，输入``帮我设计一座20米简支梁桥''，它就会给你一份完整的计算书。格式规范，内容合理，甚至附带了参考文献——虽然那些参考文献可能有一半是它编的。

\textbf{当AI能把80\%的``自重''瞬间完成，那剩下的20\%——那些AI做不了的东西——到底是什么？}

\section{效率杀不死意义，但它会让意义变哑}

2026年6月11日，师宇写了一篇博文，标题叫``有什么用呢''。他研究了AI视频生成工具——Coze、Seedance、花生AI——兴奋地做了很多尝试。然后被一句话问住了：

\begin{shiyu-inline}
``我原本兴奋于'是不是该认真做一些有趣的视频'，却被一句'有什么用呢'问住。什么算有用？必须产出、变现、转化才叫有用吗？不是所有有用的东西，都必须马上证明自己有用。有些价值正是在暂时无用时显现。''

\hfill ——师宇博文，2026年6月11日
\end{shiyu-inline}

这句话击中了问题的核心。

``这个学了有什么用''——这句话本身，就是效率思维的产物。在效率至上的世界里，一切``无用''之事都被驱逐了。诗、音乐、游戏、搭一座你很确定不会有人走的桥。这些事都没有``用''。但它们让生活值得过。

\textbf{效率杀不死意义。但效率会让意义变哑。}当所有人都只问``有什么用''，那些``没用但重要''的东西就不再被说出口。它们没有消失，它们只是沉默了。

\section{从梁桥到拱桥}

梁桥的困境，本质上是\textbf{效率与意义的矛盾}。

你想让梁桥的跨度更大？你必须让它更重。你想让它更轻？它的跨度就会变小。这是一个物理定律，没有变通的办法。

但工程师没有放弃。他们发明了拱桥。

拱桥的秘密是：\textbf{把垂直的荷载，转化为水平的推力。}压力不是被抵抗的，而是被利用的。石头不是被压垮的，而是被压紧的。\textbf{压力，创造了强度。}

这就是第二部分要讲的事。

在接下来的几章里，我们会看到：当AI冲击旧的教育系统，当学生带着ChatGPT的答案走进教室，当排名、积分、点赞成为新的课堂货币——这些压力，不是要摧毁教育的敌人。它们是\textbf{建造拱桥的石头}。

\begin{insight}
\label{insight:chap3}

\textbf{梁桥的困境，不是要我们放弃梁桥。是要我们学会造拱桥。}

效率杀不死意义。但效率会让我们忘记，意义从来不是``有用''的东西。意义是你在洗浴中心搭桥搭到关门，别人问你为什么，你只能说：``没想什么。就是塌了，想让它站住。''

那就是意义。那不需要被证明有用。
\end{insight}

\vspace{1cm}

\begin{shiyu-note}
\label{note:chap3}

\textbf{在效率的海洋里，为``无用的认真''保留一片海滩。}

有一个问题，我建议所有老师在开学第一节课问学生：

``如果这门课不算学分，没有成绩，不写入档案——你还会来上吗？''

这不是一个测试。这是一面镜子。它照出的不是学生是否勤奋，而是这门课是否值得被教。

如果答案是``不会''，那问题不在学生。在课程设计。我们花了太多时间让课程``有用''——符合大纲、覆盖考点、支撑后续课程——却忘了让课程``值得''。

值得，不是有用。值得是：即使没有任何外部奖励，你仍然愿意做这件事。就像孩子蹲在沙坑里搭三个小时的桥。就像洗浴中心里那个搭桥搭到关门的大学生。

AI时代，效率不再是稀缺品。意义才是。
\end{shiyu-note}

\vspace{1cm}

\begin{decision-point}
\label{decision:chap3}

\noindent\textbf{这一章，我们看到了效率如何杀死了意义。}

\vspace{0.3cm}

\noindent 下一次备课，做一件事：把你准备讲的内容分成两列。左边写``这个学生必须知道''，右边写``这个学生不学也可以，但我想让他们感受一下''。

\noindent 然后，把左边砍掉一半。把右边加倍。

\noindent 告诉我，那节课发生了什么。
\end{decision-point}

\vspace{0.5cm}

\begin{flushright}
{\small\color{muted} 第一部分结束。下一部分：拱桥——当AI冲击旧系统，压力如何创造强度。}
\end{flushright}
